\begin{register}{H}{pma\_cfg\regindex{n} (\regindex{n}: 0-15)} {0xBC0+0x1*\regindex{n}}
\label{regdesc:pma_cfgX}
\regfield{A}{2}{30}{{2}}%
\regfield{LOCK}{1}{29}{{0}}%
\regfield{reserved}{1}{28}{{0}}%
\regfield{ATTRIBUTE}{4}{24}{{0}}%
\regfield{reserved}{19}{5}{{0}}%
\regfield{READ}{1}{4}{{0}}%
\regfield{WRITE}{1}{3}{{0}}%
\regfield{EXECUTE}{1}{2}{{0}}%
\regfield{reserved}{1}{1}{{0}}%
\regfield{TYPE}{1}{0}{{0}}%
\reglabel{Reset}\regnewline%
\begin{regdesc}
\begin{reglist}
\label{fielddesc:A}\item [A] 配置地址类型。功能与 pmpcfg 寄存器的 A 字段相同。\\
0x0：OFF\\
0x1：TOR\\
0x2：NA4\\
0x3：NAPOT\\
(R/W)
\label{fielddesc:LOCK}\item [LOCK] 配置是否锁定对应的 pma\_cfgX 和 pma\_addrX。\\
0：不锁定\\
1：锁定，即撤销写入 pma\_cfgX 和 pma\_addrX 的权限\\
只有 CPU 复位才能解锁。\\
(R/W)
\label{fielddesc:ATTRIBUTE}\item [ATTRIBUTE] 配置要在 DRAM 属性端口上驱动的值。\\
bit 27： \\
0：可缓存 (cacheable)\\
1：不可缓存 (non-cacheable)\\
bit 26：\\
0：写回（write-back）\\
1：写直达（write-through）\\
bit 25：\\
0：写失效分配（write miss allocate）\\
1：写失效不分配（write miss no allocate）\\
bit 24：
0：读失效分配（read miss allocate）\\
1：读失效不分配（read miss no allocate）\\
(R/W)
\label{fielddesc:READ}\item [READ] 配置相应区域的读取权限。\\
0：不允许读取\\
1：允许读取\\
 (R/W)
\label{fielddesc:WRITE}\item [WRITE] 配置相应区域的写权限。\\
0：不允许写入\\
1：允许写入\\
(R/W)
\item [见下页……]
\end{reglist}\end{regdesc}
\end{register}
\addtocounter{Regfloat}{-1}
\begin{register}{H}{pma\_cfg\regindex{n} (\regindex{n}: 0-15)} {0xBC0+0x1*\regindex{n}}
\begin{regdesc}\begin{reglist}
 \item [接上页……]
\label{fielddesc:EXECUTE}\item [EXECUTE] 配置相应区域的执行权限。\\
0：不允许执行\\
1：允许执行\\
 (R/W)
\label{fielddesc:TYPE}\item [TYPE] 配置区域类型。\\
0x0：存储器区域无效（RWX 访问将被视为 0，即使配置为 1）\\
0x1：存储器区域有效（RWX 访问可用）\\
(R/W)
\end{reglist}
\end{regdesc}
\end{register}

\begin{register}{H}{pma\_addr\regindex{n} (\regindex{n}: 0-15)} {0xBD0+0x1*\regindex{n}}
\label{regdesc:pma_addrX}
\regfield{ADDR}{32}{0}{{0}}%
\reglabel{Reset}\regnewline%
\begin{regdesc}
\begin{reglist}
\label{fielddesc:PMAADDR}\item [ADDR] 配置地址。功能与 pmpaddr 寄存器相同。(R/W)
\end{reglist}
\end{regdesc}
\end{register}
